\input{preamble}

\title{Section 3.1 --- Power and Radical Functions --- Answer Key}

\NewDocumentCommand{\DomRng}{m d(] d[)}{%
$\dom #1=\eint$\\$\rng #1=%
\IfNoValueTF{#2}{%
\IfNoValueTF{#3}{%
\eint
}{
\Lint{#3}
}}{
\Rint{#2}
}$}

\newcommand{\domrng}[3]{%
$\dom #1=#2$\\$\rng #1=#3$%
}

\newlength{\lht}
\newlength{\lwid}
\newcommand{\Lboxed}[1]{%
\setlength{\lwid}{\widthof{$#1$}}
\setlength{\lht}{\heightof{ABC123}}
\begin{tikzpicture}[baseline=0]
\useasboundingbox (-0.5\lwid-3pt,0) rectangle (0.5\lwid+3pt,\lht);
\draw (0,0) node[anchor=base]{$#1$};
\draw[rounded corners=3pt] (-0.5\lwid-2pt,-3pt) rectangle (0.5\lwid+2pt,\lht+3pt);
\end{tikzpicture}%
}

\begin{document}
\maketitle
\section{Graphs of Power Functions}
Match each function below to its graph.
\begin{Nmenum}
\mitemxxxx
{$f(x)=x^2$}
{$g(x)=x^3$}
{$h(x)=x^6$}
{$k(x)=x^7$}
\end{Nmenum}
\begin{smenum}
\mitemxxxx
{%
B\LittleGraphBlank[3]{\powgraph{3.5}{3}}
}{%
D\LittleGraphBlank[3]{\powgraph{3.5}{7}}
}{%
C\LittleGraphBlank[3]{\powgraph{3.5}{6}}
}{%
A\LittleGraphBlank[3]{\powgraph{3.5}{2}}
}
\end{smenum}

\section{Domain and Range of Power Functions}
Give the domain and range of each power function.
\begin{smenum}
\mitemxxx
{\DomRng{f}}
{\DomRng{g}[6)}
{\DomRng{h}}
\mitemxxx
{\DomRng{k}(3]}
{\DomRng{p}[-4)}
{\DomRng{q}(-8]}
\end{smenum}

\section{Evaluating Roots}
Fill in the blanks to evaluate each root by trial and error. (All of these are defined.)
\begin{smenum}
\mitemxx
{$\sqrt{49}=\Lboxed{7}$ because $\Lboxed{7^2}=49$}
{$\sqrt[3]{64}=\Lboxed{4}$ because $\Lboxed{4^3}=64$}
\mitemxx
{$\sqrt[6]{729}=\Lboxed{3}$ because $\Lboxed{3^6}=729$}
{$\sqrt[5]{-32}=\Lboxed{-2}$ because $\Lboxed{(-2)^5}=-32$}
\end{smenum}
Fill in the blanks to show the two consecutive integers each root is between by trial and error.
\begin{smenum}
\mitemx
{$\Lboxed{2}<\sqrt[5]{183}<\Lboxed{3}$ because $\Lboxed{32}<183<\Lboxed{243}$}
\mitemx
{$\Lboxed{8}<\sqrt{74}<\Lboxed{9}$ because $\Lboxed{64}<74<\Lboxed{81}$}
\mitemx
{$\Lboxed{5}<\sqrt[4]{731}<\Lboxed{6}$ because $\Lboxed{625}<731<\Lboxed{1296}$}
\mitemx
{$\Lboxed{-5}<\sqrt[3]{-99}<\Lboxed{-4}$ because $\Lboxed{-125}<-99<\Lboxed{-64}$}
\end{smenum}

\section{Evaluating Roots on the Calculator}
Evaluate each root using the calculator. Round your answers to three decimal places, if necessary. If the given root is undefined, write ``undefined.''
\begin{smenum}
\mitemxxx
{$\sqrt[6]{94}\approx 2.132$}
{$\sqrt{15}\approx 3.873$}
{$\sqrt{27}\approx 5.196$}
\mitemxxx
{$\sqrt[5]{-116}\approx -2.588$}
{$\sqrt[5]{11458}\approx 6.484$}
{$\sqrt[6]{-3323}$ is undefined}
\mitemxxx
{$\sqrt[7]{-66577}\approx -4.887$}
{$\sqrt[3]{237}\approx 6.188$}
{$\sqrt[4]{-554}$ is undefined}
\mitemx
{$\sqrt[3]{202}\approx 5.867$}
\end{smenum}

\section{Domain and Range of Radical Functions}
Find the domain and range of each radical function.
\begin{smenum}
\mitemxxx
{\domrng{f}{\Lint{\frac{1}{3}}}{\Lint{3}}}
{\domrng{g}{\Rint{-2}}{\Lint{-2}}}
{\domrng{h}{\Lint{-6}}{\Rint{-5}}}
\mitemxxx
{\DomRng{k}}
{\DomRng{p}}
{\domrng{q}{\Rint{7}}{\Rint{1}}}
\end{smenum}
\end{document}
